\documentclass[border=2pt]{standalone}
\usepackage{amsmath,amssymb}
\usepackage{3dfigs}
\begin{document}

\begin{tikzpicture}[scale=1.15]
  % sphere radius R=4, point P1 at theta=20 deg, phi=10 deg
  % sphere radius R=4, point P2 at theta=15 deg, phi=-20 deg
  % sphere radius R=4, point P3 at theta=40 deg, phi=0 deg
  % ray from center
  \coordinate (O) (0,0,0) ;
  \coordinate (P1) ({4*cos(10)*sin(20)},{4*sin(10)*sin(20)},{4*cos(20)});
  \coordinate (P2) ({4*cos(-20)*sin(15)},{4*sin(-20)*sin(15)},{4*cos(15)});
  \coordinate (P3) ({4*cos(0)*sin(40)},{4*sin(0)*sin(40)},{4*cos(40)});
  \draw[bvec] (O) --($1.2*(P1)$);
  \draw[bvec] (O) --($1.2*(P2)$);
  \draw[bvec] (O) --($1.2*(P3)$);
  \tikzmath{
    \ds1 = acos(cos(-20)*cos(0)+sin(-20)*sin(0)*cos(40-15));
    \ds2 = acos(cos(10)*cos(0)+sin(10)*sin(0)*cos(40-20));
    \ds3 = acos(cos(10)*cos(-20)+sin(10)*sin(-20)*cos(20-15));
    \ssc1 = 1/sin(\ds1);
    \ssc2 = 1/sin(\ds2);
    \ssc3 = 1/sin(\ds3);
 }
  % beta2= p3-p1
  \draw[green] plot[variable=\t,domain=0:1,samples=60] ({4*cos(-10*\t)*sin(40-20*\t)},{4*sin(-10*\t)*sin(40-20*\t)},{4*cos(40-20*\t)});
  % beta3= p1-p2
  \draw plot[variable=\t,domain=0:1,samples=60] ({4*cos(10-30*\t)*sin(20-5*\t)},{4*sin(10-30*\t)*sin(20-5*\t)},{4*cos(20-5*\t)});
\end{tikzpicture}
\end{document}

  % beta1= p2-p3
  \draw[blue] plot[variable=\t,domain=0:1,samples=60] ($cos((1-\t)*\ds3)*\ssc1*(P2)+cos(\t*\ds3)*ssc1*(P3)$);;

    \draw[cdash] plot[variable=\t,domain=0:90,samples=60] 
  \draw[cdash] plot[variable=\t,domain=0:90,samples=60] (0,{4*sin(\t)},{4*cos(\t)});
  \draw[cdash] plot[variable=\t,domain=0:90,samples=60] ({4*cos(\t)},{4*sin(\t)},0);
  % meridian through P (azimuth 50 deg): pole -> P -> E
  \draw[rdash] plot[variable=\t,domain=0:90,samples=60]
       ({2.571*sin(\t)},{3.064*sin(\t)},{4*cos(\t)});
  % drop line P->Q and radius O->Q in the xy-plane
  \draw[cdash] (1.653,1.970,3.064)--(1.653,1.970,0);
  \draw[cdash] (0,0,0)--(1.653,1.970,0);
  % axes
  \draw[axis] (0,0,0)--(5,0,0) node[below left]{$x$};
  \draw[axis] (0,0,0)--(0,5.5,0) node[right]{$y$};
  \draw[axis] (0,0,0)--(0,0,5) node[above]{$z$};
  % position vector r
  \draw[rvec] (0,0,0)--(1.653,1.970,3.064);
  \node[blue!70!black] at (1.25,1.4,2.4) {$\mathbf{r}$};
  \fill (1.653,1.970,3.064) circle (1.4pt);
  % polar angle theta (z-axis to r)
  \draw[arc] plot[variable=\t,domain=0:40,samples=30]
       ({1.1*sin(\t)*0.643},{1.1*sin(\t)*0.766},{1.1*cos(\t)});
  \node at (0.42,0.5,1.45) {$\vartheta$};
  % azimuth phi (x-axis to OQ, in xy-plane)
  \draw[arc] plot[variable=\t,domain=0:50,samples=30] ({1.0*cos(\t)},{1.0*sin(\t)},0);
  \node at (1.3,0.8,0) {$\varphi$};
  % local spherical basis at P
  \draw[bvec] (1.653,1.970,3.064)--(2.190,2.610,4.060);
  \node[red!80!black,above] at (2.19,2.61,4.10) {$\mathbf{e}_r$};
  \draw[bvec] (1.653,1.970,3.064)--(2.293,2.733,2.228);
  \node[red!80!black,right] at (2.29,2.73,2.20) {$\mathbf{e}_\vartheta$};
  \draw[bvec] (1.653,1.970,3.064)--(0.657,2.806,3.064);
  \node[red!80!black,above] at (0.66,2.81,3.10) {$\mathbf{e}_\varphi$};